% !TEX root = ../../../MathModel.tex
% 负责人：罗财能。
\subsubsection{具体分析}
\label{sec:q2:analysis}

问题二将单点运输能力扩展为异构机队的多点、多架次调度。
需要同时决定每个货箱归属哪个架次、沿途访问顺序、运输机型、具体机体、电池及准备开始时刻。
一个架次允许访问多个服务区，每次投送后载荷减少，后续航段的能耗应随之重算；各站均须完成下降和交接后重新爬升，不能把多点路线仅视为水平距离的合并。

调度的关键在于机体与电池具有不同的释放时刻。
运输机返航后可换用另一组满电电池继续任务，而卸下的电池需要按返航SOC充至100\%才可复用。
因此，实体机数量充足并不意味着可以立即出发，单独计算机体占用会低估电池周转造成的等待。
同型电池允许在不同运输机间共享，跨机型则不能混用。

本问还须区分两类时间要求：医疗物资的期望送达时间和首批保障截止时间属于硬约束，普通物资的期望时间用于评价配送及时性。
本文优先降低加权延迟，在其相同的方案中依次比较全部运输机返航完成时间、运输能耗和架次数。
该顺序体现抢救时效优先的选择，不把四个指标混为同一单位，也不假定它们能同时达到最小。

第二题从原始运输数据独立求解，沿用公共地形、载荷、能耗和充电规则；不施加通信或中继资源约束，也不使用第三题的已求解方案作为输入。
采用有限候选构造与固定结构排程改进相结合的方法，最后独立核验实际浮点时序，避免仅凭优化器状态判断资源可行性。
